% Source: https://artofproblemsolving.com/community/c3h2934055_2016_emcc_guts_round_p13_p24__exeter_math_club_competition
% Stable identity: aops-topic-2934055-q15
Put $C=(0,0),A=(1,0),B=(0,1)$. Since $CDE$ is right isosceles, $D=(0,t),E=(t,0)$ with $0<t<1$. For $AEF$ to be right isosceles, either its right angle is at $E$, giving $F=(t,1-t)$, or at $F$, giving $F=(\dfrac{1+t}{2},\dfrac{1-t}{2})$; a right angle at $A$ cannot put $F$ on $AB$. Similarly $BFD$ gives $F=(1-t,t)$ or $F=(\dfrac{1-t}{2},\dfrac{1+t}{2})$. Equating these possibilities gives $t=\dfrac{1}{2}$ or $t=\dfrac{1}{3}$; the fourth pairing gives the forbidden $t=0$. In both allowed cases direct dot products verify $DEF$ is right isosceles. Thus the distinct lengths are $\dfrac{\sqrt2}{2},\dfrac{\sqrt2}{3}$, summing to $\boxed{5\dfrac{\sqrt2}{6}}$.
\begin{asy}
usepackage("amsmath");
size(7cm);
for(int k=0;k<2;++k){pair off=(1.5*k,0);real t=k==0?.5:1/3;pair C=off,A=off+(1,0),B=off+(0,1),D=off+(0,t),Ep=off+(t,0),F=off+(t,1-t);draw(A--B--C--cycle);draw(D--Ep--F--cycle,blue);label("$A$",A,S);label("$B$",B,N);label("$C$",C,SW);label("$D$",D,W);label("$E$",Ep,S);label("$F$",F,E);}
\end{asy}
